% \usepackage{palatino} % Palatino
% \usepackage{mathpazo} % Palatino
\usepackage{mathptmx} % Times
% \renewcommand{\ttdefault}{lmtt}

% For badges
\usepackage{tikz}
\usepackage{xcolor}

% Define badge colors (pastel, matching website)
\definecolor{utdcolor}{RGB}{173, 216, 230}     % Pastel blue for UTD24
\definecolor{ftcolor}{RGB}{255, 182, 193}      % Pastel red for FT50
\definecolor{abscolor}{RGB}{255, 218, 185}     % Pastel orange for ABS
\definecolor{awardcolor}{RGB}{180, 216, 180}   % Pastel green for awards/TOP CS (muted sage)
% Text colors for badges (dark for readability on pastel backgrounds)
\definecolor{utdtext}{RGB}{0, 70, 100}         % Dark blue text
\definecolor{fttext}{RGB}{139, 0, 50}          % Dark red text
\definecolor{abstext}{RGB}{180, 100, 0}        % Dark orange text
\definecolor{awardtext}{RGB}{0, 100, 0}        % Dark green text

% Badge command - creates small rounded box (middle-aligned)
\newcommand{\badge}[3]{%
  \tikz[baseline=-.6ex]{%
    \node[fill=#1, text=#2, rounded corners=2pt, inner sep=2pt, font=\tiny\bfseries] {#3};%
  }%
}

% Specific badge commands for journal rankings
\newcommand{\utd}{\badge{utdcolor}{utdtext}{UTD24}}
\newcommand{\ft}{\badge{ftcolor}{fttext}{FT50}}
\newcommand{\absfourstar}{\badge{abscolor}{abstext}{ABS4*}}
\newcommand{\absfour}{\badge{abscolor}{abstext}{ABS4}}
\newcommand{\absthree}{\badge{abscolor}{abstext}{ABS3}}
\newcommand{\absgen}{\badge{abscolor}{abstext}{ABS}}

% TOP CS badge
\newcommand{\topcs}{\badge{awardcolor}{awardtext}{TOP CS}}

% Award badges
\newcommand{\bestpaper}{\badge{awardcolor}{awardtext}{Best Paper}}
\newcommand{\bestposter}{\badge{awardcolor}{awardtext}{Best Poster}}
\newcommand{\bestpaperfinalist}{\badge{awardcolor}{awardtext}{Best Paper Finalist}}
\newcommand{\beststudentpaper}{\badge{awardcolor}{awardtext}{Best Student Paper}}

% Publication entry with hanging indentation
% Usage: \pubentry{[J39]}{Content here...}
\newcommand{\pubentry}[2]{%
  \noindent\hangindent=2.8em\hangafter=1%
  \makebox[2.5em][r]{\textbf{#1}}~#2\par\vspace{0.5em}%
}

\usepackage{lastpage}
\usepackage{fancyhdr}
\usepackage{ifthen}
\pagestyle{fancy}
\fancyhf{}
\renewcommand\headrulewidth{0pt}
\lfoot{\ifthenelse{\value{page}=1}{}{\small Hyunwoo Park (Last Updated: \today)}}
\rfoot{\ifthenelse{\value{page}=1}{}{\small Page \thepage\ of \pageref*{LastPage}}}

% \usepackage{titlesec}
% \titlespacing*{\section}
% {0pt}{5.5ex plus 1ex minus .2ex}{4.3ex plus .2ex}
% \titlespacing*{\subsection}
% {0pt}{5.5ex plus 1ex minus .2ex}{4.3ex plus .2ex}

\setlength{\parskip}{0pt}
% \setlength{\parsep}{10pt}
% \setlength{\headsep}{0pt}
% \setlength{\topskip}{0pt}
% \setlength{\topmargin}{0pt}
% \setlength{\topsep}{0pt}
% \setlength{\partopsep}{0pt}

\makeatletter
\let\origsection\section
\renewcommand\section{\@ifstar{\starsection}{\nostarsection}}

\newcommand\nostarsection[1]
{\sectionprelude\origsection{#1}\sectionpostlude}

\newcommand\starsection[1]
{\sectionprelude\origsection*{#1}\sectionpostlude}

\newcommand\sectionprelude{%
  \vspace{1em}
}

\newcommand\sectionpostlude{%
  \vspace{0em}
}
\makeatother


% \itemsep10pt\parskip0pt\parsep0pt
\setlength{\itemsep}{10pt}


\usepackage{etaremune}
\usepackage{fontawesome}
\usepackage{etoolbox}

% Toggle for reverse-numbered lists
\newbool{reversenum}
\boolfalse{reversenum}

\let\oldenumerate\enumerate
\let\oldendenumerate\endenumerate

\renewenvironment{enumerate}{%
  \ifbool{reversenum}{\etaremune}{\oldenumerate}%
}{%
  \ifbool{reversenum}{\endetaremune}{\oldendenumerate}%
}

% Override Pandoc's label definition after \tightlist
\let\oldtightlist\tightlist
\renewcommand{\tightlist}{%
  \oldtightlist
  \ifbool{reversenum}{\def\labelenumi{[\theenumi]}}{}%
}

\newcommand{\reversenumstart}{\booltrue{reversenum}}
\newcommand{\reversenumstop}{\boolfalse{reversenum}}

